% Source: https://artofproblemsolving.com/community/c3h3558367_can_someone_explain_this_one
% Stable identity: aops-topic-3558367-q01
The identity $\sum_{k=1}^n k(k+1)=\dfrac{n(n+1)(n+2)}{3}$ follows by adding the usual sums of squares and integers. Thus the required sum is $\dfrac{n(n+1)(n+2)}{2}$. At $n=14$ it is 1680; at $n=15$ it is 2040. Positive terms make the sums strictly increasing, so $n=14$ and the final term is $\dfrac{3\cdot 14\cdot 15}{2}=\boxed{315}$.
